\documentclass[11pt]{article}
\usepackage{amsmath,amssymb,amsthm}
\setlength{\textwidth}{6.3in}
\setlength{\oddsidemargin}{0.1in}
\setlength{\evensidemargin}{0.1in}
\setlength{\textheight}{8.6in}
\setlength{\topmargin}{-0.3in}
\setlength{\headheight}{14pt}
\setlength{\headsep}{22pt}
\setlength{\parindent}{1.2em}
\setlength{\parskip}{3pt}
\newcommand{\R}{\mathbb R}
\newcommand{\N}{\mathbb N}
\newcommand{\Nzero}{\mathbb N_0}
\newcommand{\gra}{\operatorname{gra}}
\newcommand{\dom}{\operatorname{dom}}
\newcommand{\sgn}{\operatorname{sgn}}
\newcommand{\pair}[2]{\langle #1,#2\rangle}
\newtheorem{theorem}{Theorem}
\newtheorem{proposition}{Proposition}
\newtheorem{lemma}{Lemma}
\newtheorem{example}{Example}
\newtheorem{corollary}{Corollary}
\theoremstyle{definition}
\newtheorem{definition}{Definition}
\newtheorem{assumption}{Assumption}
\theoremstyle{remark}
\usepackage{url}
\newtheorem{remark}{Remark}
\newcommand{\dist}{\operatorname{dist}}
\pagestyle{plain}

\makeatletter
\def\bstctlcite{\@ifnextchar[{\@bstctlcite}{\@bstctlcite[@auxout]}}
\def\@bstctlcite[#1]#2{\@bsphack
  \@for\@citeb:=#2\do{%
    \edef\@citeb{\expandafter\@firstofone\@citeb}%
    \if@filesw
      \immediate\write\csname #1\endcsname{\string\citation{\@citeb}}%
    \fi}%
  \@esphack}
\makeatother


\title{Nonmaximal sums of maximally monotone operators under Rockafellar's constraint qualification}

\author{Weifeng Yang\thanks{\texttt{Email:ywf841673182@gmail.com}}}
% Yunnan Earthquake Agency, Kunming, China

\date{}
% September 7, 2026

% \begin{document}



\begin{document}
\bstctlcite{IEEEexample:BSTcontrol}
\maketitle
% \footnotetext{\texttt{Email:ywf841673182@gmail.com}}


\begin{abstract}


We construct counterexamples to Rockafellar's sum conjecture in which two maximally monotone operators satisfy the interior-domain condition but their sum is not maximally monotone, thereby providing the first complete disproof of the conjecture. 
We establish a general construction theorem that computes the entire monotone polar of a class of graphs, gives a necessary and sufficient condition for their maximal monotonicity, and shows how a positive rank-one perturbation yields a nonmaximal sum under this condition. We verify the theorem's hypotheses and its maximality criterion on $c_0$, thereby obtaining a counterexample to the conjecture. Furthermore, we construct a bounded linear surjection from $\ell^1$ onto $c_0$ and use it to obtain the counterexample on $\ell^1$. Lean formalizations of the $c_0$ counterexample and the pullback lemma are also provided. 





\end{abstract}


\noindent\textbf{Keywords:} maximally monotone operator, sum theorem,
nonreflexive Banach space, monotone polar, rank-one operator.



%%介绍部分还得改下，介绍一下为解决这个问题而诞生的思想与工具

\section{Introduction}
\label{sec:introduction}
In this paper, we disprove Rockafellar's sum conjecture. 
Let $X$ be a real Banach space, $X^*$ be its continuous dual, and $A,B:X\rightrightarrows X^*$ be maximally monotone operators. Their pointwise sum is defined by $(A+B)x=\{a+b:a\in Ax,\ b\in Bx\}$. This pointwise sum is monotone, but maximal monotonicity is not automatic.  
In 1970, Rockafellar \cite{rockafellar1970} proved that $A+B$ is maximally monotone when $X$ is reflexive and
\begin{equation}
\dom A\cap\operatorname{int}(\dom B)\ne\varnothing,
\label{eq:original-cq}
\end{equation}
where the interior is taken in the norm topology of $X$. 
The question of whether the reflexivity assumption can be removed while retaining the classical constraint qualification in Eq. (\ref{eq:original-cq}) is known as Rockafellar's sum problem, or Rockafellar's sum conjecture. This conjecture has remained open for more than five decades and has been described as both the most important and the most famous open problem in monotone operator theory \cite{yao2010,borweinyao2012}. 



% Rockafellar's sum conjecture asks whether Eq. (\ref{eq:original-cq}) still guarantees that $A+B$ is maximally monotone when $X$ is an arbitrary real Banach space. 






To address Rockafellar's sum conjecture beyond reflexive spaces, a major line of research is to identify additional structural assumptions on the operators under which the sum theorem remains valid. 
For example, Yao \cite{yao2010} proves maximality when $A$ is of type~(FPV) and $B$ has full domain. Borwein and Yao \cite{borweinyao2012} prove maximality under Eq. (\ref{eq:original-cq}) when $A$ is a linear relation. Voisei \cite{voisei2010} treats the sum of an operator of type~(FPV) with convex domain and a normal-cone operator. Verona and Verona \cite{veronaverona2022} prove the sum theorem under Eq. (\ref{eq:original-cq}) when $A$ is $\mathcal C_0$-maximal, meaning that  $A+N_C$ is maximally monotone for every closed convex set $C$ with $\dom A\cap\operatorname{int}C\ne\varnothing$, where $N_C$ denotes the normal-cone operator of $C$. 
Furthermore, Bo\c{t}, Bueno, and Simons \cite{botbuenosimons2016} recall that a positive answer to Rockafellar's sum problem would imply that every maximally monotone operator is of type~(FPV). 




A complementary line of research approaches this sum conjecture through convex representations of monotone operators, allowing maximal monotonicity and sum theorems to be studied by using convex analysis. First, Fitzpatrick \cite{fitzpatrick1988} represents maximally monotone operators by convex functions on $X\times X^*$, and Burachik and Svaiter \cite{burachiksvaiter2002} study the associated family of convex representatives through enlargements. Then, within this framework, Marques Alves and Svaiter \cite{marquesalvessvaiter2008br} give sufficient conditions on a convex function and its conjugate for the represented operator to be maximally monotone in a nonreflexive Banach space. Voisei and Z\u{a}linescu \cite{voiseizalinescu2009} study the resulting class of strongly representable operators and develop its calculus rules. Using convex representations, Penot \cite{penot2004} obtains results for compositions and sums in reflexive spaces, and Marques Alves and Svaiter \cite{marquessvaiter2012} establish a sum theorem in general Banach spaces under a qualification condition involving the domains of Fitzpatrick representations. However, despite these advances, Rockafellar's sum conjecture remained unresolved.







% Sufficient conditions for maximality of sums have also been obtained through convex representations of monotone operators. Fitzpatrick \cite{fitzpatrick1988} associates a convex function with a monotone graph, and Burachik and Svaiter \cite{burachiksvaiter2002} study the associated family of convex representatives through enlargements. Using convex representations, Penot \cite{penot2004} develops results for compositions and sums in reflexive spaces, and Marques Alves and Svaiter \cite{marquessvaiter2012} prove a sum theorem in general Banach spaces under a qualification on the domains of Fitzpatrick representations. Lower bounds by the duality products for a proper lower semicontinuous convex function and its conjugate yield a maximally monotone operator \cite{marquesalvessvaiter2008br}. Voisei and Z\u{a}linescu \cite{voiseizalinescu2009} study the resulting class and its calculus.

In this paper, we give a general theorem for constructing an operation to obtain the first complete disproof of Rockafellar's sum conjecture, thereby resolving the problem after more than five decades. 
Specifically, we develop a general construction theorem that computes the entire monotone polar of a class of graphs, characterizes their maximality, and identifies a positive rank-one perturbation whose addition makes the sum nonmaximal. To satisfy the theorem's assumptions on $c_0$, we couple copies of the triangular map through a Lipschitz curve, and then obtain a counterexample to the conjecture. A bounded linear surjection then transfers this counterexample to standard $\ell^1$.

\subsection{Contributions}
The contributions of this paper are as follows.

(1) We establish a general construction theorem for counterexamples to Rockafellar's sum conjecture (Theorem~\ref{thm:endpoint-criterion}). Under Assumption~\ref{ass:construction}, the theorem computes the entire monotone polar of the graph in Eq. (\ref{eq:abs-full-graph}) and characterizes its maximality by the criterion in Eq. (\ref{eq:abs-endpoint-exclusion}). When this criterion holds, the graph defines a maximally monotone operator whose sum with a specified everywhere-defined positive rank-one operator is not maximally monotone, although the two operators satisfy Eq. (\ref{eq:original-cq}).

(2) We construct an explicit counterexample on $c_0$ (Theorem~\ref{thm:C}) by verifying the assumptions of this theorem. The second operator is bounded, positive, rank-one, and defined everywhere, and the first operator has nonempty domain. Thus the original interior-domain condition holds. We also establish a finite radial bound at the origin for the first operator (Definition~\ref{def:radial-bound}).

(3) We construct a counterexample on standard $\ell^1$ (Theorem~\ref{thm:L}) by pulling back the first operator through a specified bounded linear surjection onto $c_0$. By applying Lemma~\ref{lem:surjection-transport}, we prove that both operators are maximally monotone in $\ell^1\times\ell^\infty$. We establish nonmaximality of their sum by showing that $(0,0)$ lies outside its graph and is monotonically related to every point of that graph. 

The paper is organized as follows. Section~\ref{sec:definitions} introduces the notation and basic definitions. Section~\ref{sec:endpoint-criterion} proves the structural theorem and the pullback lemma. Section~\ref{sec:counterexamples} constructs the counterexamples on $c_0$ and standard $\ell^1$ and derives their stated consequences. Section~\ref{sec:conclusion} concludes the paper.


\section{Preliminaries}
\label{sec:definitions}
Let $\N=\{1,2,\ldots\}$ and $\N_0=\{0,1,\ldots\}$. Let $X$ be a real Banach space, we denote its continuous dual by $X^*$. We define $\pair{x}{a}=a(x)$ for $x\in X$ and $a\in X^*$. For a real Hilbert space $H$, its inner product is denoted by $(u,v)_H$. We use the norm topology for interiors and closures unless stated otherwise. 

For a graph $G\subset X\times X^*$, we also define $\dom G=\{x:\exists a,\ (x,a)\in G\}$ and $\operatorname{ran}G=\{a:\exists x,\ (x,a)\in G\}$. 
For a bounded linear map $P:X\to X^*$, we call $P$ positive if $\pair{x}{Px}\ge0$ for every $x\in X$, and rank one if $\dim(\operatorname{ran}P)=1$. 

Next, we introduce monotonicity and the monotone polar  \cite{martinezsvaiter2005,bauschkemclarensendov2006} as follows. 

\begin{definition}[Monotonicity and the monotone polar]
\label{D1}
For a graph $G\subset X\times X^*$, define
\begin{equation}
G^\mu=\{(z,p)\in X\times X^*:
 \pair{z-x}{p-a}\ge0\quad\forall(x,a)\in G\}.
\label{eq:D1}
\end{equation}
A point $(z,p)$ is monotonically related to $G$ if $(z,p)\in G^\mu$.
The graph is monotone when $\pair{x-y}{a-b}\ge0$ for all $(x,a),(y,b)\in G$, and the graph is maximally monotone if it is monotone and there is no monotone graph $\widetilde G\subset X\times X^*$ such that $G\subsetneq\widetilde G$. 
\end{definition}

By Eq. (\ref{eq:D1}), $(0,0)\in G^\mu$ if and only if $\pair{x}{a}\ge0$ for every $(x,a)\in G$. Furthermore, from Definition \ref{D1}, we can prove the following proposition, which shows that maximal monotonicity can be characterized directly by the monotone polar. 


\begin{proposition}
A monotone graph $G\subset X\times X^*$ is maximally monotone if and only if $G=G^\mu$.
\end{proposition}
\begin{proof}
A point in $G^\mu\setminus G$ can be adjoined to $G$, and every point of a monotone extension belongs to $G^\mu$.
\end{proof}


Next, we introduce a quantity that determines whether there exists a constant $C\ge0$ such that $\pair{x}{a}\ge-C\|x\|$ for every $(x,a)\in G$. 

\begin{definition}[Radial bound at the origin]
\label{def:radial-bound}
For a nonempty graph $G$ with $0\notin\dom G$, define
\begin{equation}
\mathcal V(G)=\sup_{(x,a)\in G}
 \frac{\max\{0,-\pair{x}{a}\}}{\|x\|}.
\label{eq:D3}
\end{equation}
We say that $G$ has a finite radial bound at the origin if $\mathcal V(G)<\infty$.
\end{definition}
For every $C\ge0$, Eq. (\ref{eq:D3}) gives that 
\[
\mathcal V(G)\le C
\quad\Longleftrightarrow\quad
\pair{x}{a}\ge-C\|x\|
\quad\text{for all }(x,a)\in G.
\]
This means $\mathcal V(G)<\infty$ if and only if there exists $C\ge0$ such that $\pair{x}{a}\ge-C\|x\|$ for every $(x,a)\in G$, and $\mathcal V(G)$ is the smallest such constant in this case. 





\section{Maximal monotonicity and nonmaximal sums}
\label{sec:endpoint-criterion}

In this section, we establish a general construction theorem for maximally monotone operators whose sums with a positive rank-one operator with full domain are not maximally monotone. We also prove a pullback lemma showing that maximal monotonicity is preserved under pullback by a bounded linear surjection.  






We first introduce the graph used in our construction. The construction consists of a Hilbert-valued profile, a scalar parameter, and a linear map from the dual space to the primal space. 


\begin{definition}[Hilbert spaces and norms]
\label{def:hilbert-spaces-norms}
Let $H_0$ be a real Hilbert space with inner product $(\cdot,\cdot)_{H_0}$ and induced norm
\[
\|u\|_{H_0}=\sqrt{(u,u)_{H_0}}\qquad(u\in H_0).
\]
Define $H=\R\oplus H_0$ to be $\R\times H_0$ equipped with the inner product
\[
((s,u),(t,v))_H=st+(u,v)_{H_0}
\qquad(s,t\in\R,\ u,v\in H_0).
\]
Its induced norm is
\[
\|(s,u)\|_H=\sqrt{s^2+\|u\|_{H_0}^2}.
\]
\end{definition}


\begin{definition}[Auxiliary structure]
\label{def:construction-maps}
Let $E:X^*\to X$ and $M:X^*\to H$ be bounded linear maps, and let
$g\in X^*\setminus\{0\}$, we define 
\[
h=Eg
\]
and define the scalar functional
\[
r(a)=\pair{h}{a}.
\]
We further define
\[
La=-Ea+r(a)h,
\qquad
K=\ker M\cap\ker r.
\]
\end{definition}

% We use $r(a)$ as a scalar parameter and constrain $Ma$ to lie on a prescribed curve at this parameter, thereby obtaining the graph used in our construction. 


We use $r(a)$ as a scalar parameter and constrain $Ma$ to equal the value of a prescribed curve at $r(a)$, thereby obtaining the graph used in our construction.

% We now prescribe the relation between these two quantities.

\begin{definition}[The constrained graph]
\label{def:construction-graph}
Let $\omega:(0,\infty)\to H_0$ be $\ell$-Lipschitz, where $0\le\ell\le1$. We define 
\[
\begin{aligned}
J&=(-\infty,-1)\cup(1,\infty),\\
F(t)&=(\sgn t,\omega(|t|-1))\qquad(t\in J), 
\end{aligned}
\]
where \[
\operatorname{sgn}(t)=
\begin{cases}
1, & t>0,\\
0, & t=0,\\
-1, & t<0.
\end{cases}
\] 
The admissible dual variables and the graph are
\begin{equation}
\begin{aligned}
D&=\{a\in X^*:r(a)\in J,\ Ma=F(r(a))\},\\
G&=\{(La,a):a\in D\}.
\end{aligned}
\label{eq:abs-full-graph}
\end{equation}
\end{definition}




Next, we introduce the following assumptions, which provide the structure needed to compute the monotone polar $G^\mu$ explicitly and characterize the maximality of $G$. 



\begin{assumption}\label{ass:construction}
For the maps and graph in Definitions~\ref{def:construction-maps} and~\ref{def:construction-graph}, assume:
\begin{enumerate}
\renewcommand{\theenumi}{H\arabic{enumi}}
\renewcommand{\labelenumi}{(\theenumi)}
\item\label{ass:bilinear} $Mg=0$, and
\begin{equation}
\pair{Ea}{b}+\pair{Eb}{a}=2(Ma,Mb)_H
\qquad(a,b\in X^*).
\label{eq:abs-bilinear}
\end{equation}

\item\label{ass:annihilator} The subspace $K$ satisfies
\begin{equation}
\{x\in X:\pair{x}{k}=0\ \text{for all }k\in K\}=\R h.
\label{eq:abs-annihilator}
\end{equation}

\item\label{ass:curve-limit} There exists $\eta\in H_0$ such that
\begin{equation}
\begin{gathered}
\omega(s)\longrightarrow\eta\text{ in }H_0
\quad\text{as }s\downarrow0,\\
0<S:=\|\eta\|_{H_0}^2<1,\\
\|\omega(s)\|_{H_0}^2\le S\qquad(s>0).
\end{gathered}
\label{eq:abs-tail}
\end{equation}

\item\label{ass:realization} For every $t\in J$, there exists $a^t\in X^*$
satisfying $r(a^t)=t$ and $Ma^t=F(t)$.
\end{enumerate}
\end{assumption}
% By Assumption~\ref{ass:construction}\textup{(\ref{ass:realization})}, each parameter fibre
% $\{a\in X^*:r(a)=t,\ Ma=F(t)\}$ is nonempty. Since $K=\ker M\cap\ker r$, this fibre is exactly $a^t+K$. Together with the other conditions in Assumption~\ref{ass:construction}, these affine fibres reduce the polar inequalities to the scalar parameter $r(a)$ and the Hilbert-valued profile $Ma$.




Therefore, under Assumption~\ref{ass:construction}, we establish the following theorem, which gives a general construction mechanism for counterexamples satisfying the interior-domain condition in Eq. (\ref{eq:original-cq}). Specifically, this theorem computes the entire monotone polar $G^\mu$ and characterizes the maximal monotonicity of $G$ by the condition that no $(p,\tau)\in X^*\times[-1,1]$ satisfies $r(p)=\tau$ and $Mp=(\tau,\eta)$.  This system characterizes exactly the points $(Lp,p)$ outside $G$ that can be adjoined to $G$ while preserving monotonicity. Consequently, when this system has no solution, the operator with graph $G$ is maximally monotone. By combining this maximality criterion with the positive rank-one perturbation $\mathcal P x=\pair{x}{g}g$, this theorem provides a general mechanism for constructing pairs of maximally monotone operators that satisfy Eq. (\ref{eq:original-cq}) while their sum is not maximally monotone. 









\begin{theorem}[Construction of nonmaximal sums]
\label{thm:endpoint-criterion}
For the graph $G$ in Definition~\ref{def:construction-graph}, suppose Assumption~\ref{ass:construction} holds and define
\[
\widehat F(t)=
\begin{cases}
F(t),&|t|>1,\\
(t,\eta),&|t|\le1.
\end{cases}
\]
Then the following statements hold.
\begin{enumerate}
\renewcommand{\theenumi}{\roman{enumi}}
\renewcommand{\labelenumi}{(\theenumi)}
\item\label{thm:polar-part} The entire monotone polar of $G$ in $X\times X^*$ is
\begin{equation}
G^\mu=\{(Lp,p):p\in X^*,\ Mp=\widehat F(r(p))\}.
\label{eq:abs-entire-polar}
\end{equation}
Moreover, this graph is the unique maximally monotone extension of $G$.

\item\label{thm:maximality-part} The graph $G$ is maximally monotone if and only if
\begin{equation}
\nexists(p,\tau)\in X^*\times[-1,1]:
\quad r(p)=\tau,\quad Mp=(\tau,\eta).
\label{eq:abs-endpoint-exclusion}
\end{equation}

\item\label{thm:sum-part} Suppose Eq. (\ref{eq:abs-endpoint-exclusion}) holds.
Let $\mathcal A$ be the operator with graph $G$, and define $\mathcal P x=\pair{x}{g}g$. Then
\[
0\notin\dom\mathcal A,\qquad \mathcal V(G)\le S\|g\|.
\]
Both $\mathcal A$ and $\mathcal P$ are maximally monotone and satisfy
\[
\dom\mathcal A\cap\operatorname{int}(\dom\mathcal P)\ne\varnothing.
\]
Their sum is not maximally monotone. More precisely,
\[
(0,0)\in
\bigl(\gra(\mathcal A+\mathcal P)\bigr)^\mu
\setminus\gra(\mathcal A+\mathcal P).
\]
\end{enumerate}
\end{theorem}
\begin{proof}
\noindent\textit{\textup{(\ref{thm:polar-part})}}
Let $a=b=g$ in Eq. (\ref{eq:abs-bilinear}), from $Mg=0$, we know that $r(g)=0$, thus we have 
\begin{equation}
\begin{aligned}
\pair{Ea}{g}&=-r(a),\\
\pair{La}{g}&=r(a),\\
\pair{La-Lb}{a-b}&=|r(a)-r(b)|^2-\|Ma-Mb\|^2.
\end{aligned}
\label{eq:abs-pairing}
\end{equation}
Taking $t=2$ in Assumption~\ref{ass:construction}\textup{(\ref{ass:realization})}, we can obtain $r(a^2)=2$, thus $h\ne0$.

Next, to establish the monotonicity properties needed to compute $G^\mu$, we first prove that $\widehat F$ is $1$-Lipschitz on $\R$. 
We extend $\omega$ to $[0,\infty)$ by setting $\omega(0)=\eta$, then let $\chi(t)=\max\{-1,\min\{1,t\}\}$ and $\zeta(t)=\max\{|t|-1,0\}$, thus, when $t$ and $u$ have the same sign, we have 
\[
|\chi(t)-\chi(u)|+|\zeta(t)-\zeta(u)|=|t-u|.
\]
For $t\ge0\ge u$, we obtain
\[
\begin{aligned}
|\chi(t)-\chi(u)|+|\zeta(t)-\zeta(u)|
&\le \chi(t)-\chi(u)+\zeta(t)+\zeta(u)\\
&=t-u.
\end{aligned}
\]
Interchanging $t,u$ covers the remaining case, thus
\[
|\chi(t)-\chi(u)|+|\zeta(t)-\zeta(u)|\le|t-u|.
\]
From the above formulations and $\widehat F(t)=(\chi(t),\omega(\zeta(t)))$, we can obtain 
\[
\|\widehat F(t)-\widehat F(u)\|^2
\le|\chi(t)-\chi(u)|^2+\ell^2|\zeta(t)-\zeta(u)|^2
\le|t-u|^2.
\]
Therefore, $\widehat F$ is $1$-Lipschitz on $\R$. Let 
\[
\widetilde G
=
\{(Lp,p):p\in X^*,\ Mp=\widehat F(r(p))\}.
\]
For any $(Lp,p),(Lq,q)\in\widetilde G$, Eq. (\ref{eq:abs-pairing}) and the $1$-Lipschitz continuity of $\widehat F$ give
\[
\pair{Lp-Lq}{p-q}
=
|r(p)-r(q)|^2-\|Mp-Mq\|^2
\ge0.
\]
Thus $\widetilde G$ is monotone. Moreover, $G\subseteq\widetilde G$ since $Ma=F(r(a))=\widehat F(r(a))$ for every $a\in D$, we know that every point of $\widetilde G$ is monotonically related to $G$, this implies
\[
\widetilde G\subseteq G^\mu.
\]


Next, we prove the reverse inclusion. Let $(x,p)\in G^\mu$, for $a\in D$ and $t=r(a)$, Eq. (\ref{eq:abs-bilinear}) gives
\begin{align}
\pair{x-La}{p-a}
={}&\pair{x}{p}-\pair{x+Ep}{a}
       +2(Ma,Mp)_H-t r(p)+t^2-\|Ma\|^2. \label{eq1}
\end{align}
For $k\in K$ and $\theta\in\R$, we have
$r(a+\theta k)=r(a)$ and $M(a+\theta k)=Ma$, thus $a+\theta k\in D$.
Therefore, from $(x,p)\in G^\mu$, substituting $a+\theta k$ into Eq. (\ref{eq1}), we have
\[
0\le
\pair{x-La}{p-a}
-\theta\pair{x+Ep}{k}
\qquad(\theta\in\R).
\]
This implies $\pair{x+Ep}{k}=0$ for every $k\in K$.
Consequently, from Eq. (\ref{eq:abs-annihilator}), we have
\begin{center}
$x=-Ep+vh$ for some $v\in\R$.    
\end{center}



We now rewrite the condition $(x,p)\in G^\mu$ in terms of the scalar parameter $t$. 
By Eq. (\ref{eq:abs-bilinear}) with $a=b=p$, we have $\pair{Ep}{p}=\|Mp\|_H^2$.  Therefore, define $\tau=\frac{v+r(p)}2,~\qquad \nu=\frac{v-r(p)}2$, from $x=-Ep+vh$, $Ma=F(t)$, and $\pair{h}{a}=t$, Eq. (\ref{eq1}) becomes 
\begin{equation}
\pair{x-La}{p-a}
=(t-\tau)^2-\nu^2-\|F(t)-Mp\|_H^2.
\label{eq:abs-polar-square}
\end{equation}
Since $(x,p)\in G^\mu$, we have 
\[
\pair{x-La}{p-a}\ge0
\qquad\text{for every }a\in D.
\]
Assumption~\ref{ass:construction}\textup{(\ref{ass:realization})} ensures that every $t\in J$ is realized by some $a^t\in D$. Therefore, from Eq. (\ref{eq:abs-polar-square}) and the above formulation, we have 
\begin{equation}
(t-\tau)^2-\nu^2-\|F(t)-Mp\|_H^2\ge0
\qquad\text{for every }t\in J.
\label{eq:abs-polar-square1}
\end{equation}
Next, we distinguish the cases $|\tau|>1$ and $|\tau|\le1$ in Eq. (\ref{eq:abs-polar-square1}).





\emph{Case 1: $|\tau|>1$.}  Since $\tau\in J$, we can take $t=\tau$ in Eq. (\ref{eq:abs-polar-square1}), which gives $\nu=0$ and $Mp=F(\tau)$. Thus $v=r(p)=\tau$ and $x=Lp$.

\emph{Case 2: $|\tau|\le1$.} In Eq. (\ref{eq:abs-polar-square1}), let $Mp=(b,z)\in H$, and $t\to1^+$ as well as $t\to-1^-$, thus from Assumption~\ref{ass:construction}\textup{(\ref{ass:curve-limit})}, we have 
\[
\begin{aligned}
(b-1)^2+\|z-\eta\|^2+\nu^2&\le(1-\tau)^2,\\
(b+1)^2+\|z-\eta\|^2+\nu^2&\le(1+\tau)^2.
\end{aligned}
\]
Multiplying these two inequalities by the nonnegative weights $\frac{1+\tau}{2}$ and $\frac{1-\tau}{2}$, and adding these two inequalities, we can obtain
\[
(b-\tau)^2+\|z-\eta\|^2+\nu^2\le0.
\]
This obviously means $b=\tau$, $z=\eta$, and $\nu=0$, thus $Mp=(\tau,\eta)$, $r(p)=v=\tau$, and $x=Lp$. 
In both cases, $x=Lp$ and $Mp=\widehat F(r(p))$, this clearly means $(x,p)\in\widetilde G$, thereby proving $G^\mu\subseteq\widetilde G$. Therefore, together with $\widetilde G\subseteq G^\mu$ proved above, Eq. (\ref{eq:abs-entire-polar}) is true.  





Moreover, since $\widetilde G$ is monotone and Eq. (\ref{eq:abs-entire-polar}) gives $G^\mu=\widetilde G$, the graph $G^\mu$ is monotone. Since $G^\mu$ is monotone and $G\subseteq G^\mu$, we know that $G^\mu\subseteq (G^\mu)^\mu\subseteq G^\mu$, which means $(G^\mu)^\mu=G^\mu$, thus $G^\mu$ is maximally monotone. 
If $\widehat G$ is any monotone extension of $G$, then 
\[
G\subseteq\widehat G\subseteq G^\mu.
\]
This means that if $\widehat G$ is maximally monotone, the monotonicity of $G^\mu$ implies $\widehat G=G^\mu$. Therefore, $G^\mu$ is the unique maximally monotone extension of $G$.  







\noindent\textit{\textup{(\ref{thm:maximality-part})}}
Since $G$ is monotone, $G$ is maximally monotone if and only if $G=G^\mu$. Therefore, by Eq. (\ref{eq:abs-entire-polar}) and Definition~\ref{def:construction-graph}, we have 
\[
G^\mu\setminus G
=
\{(Lp,p):p\in X^*,\ r(p)\in[-1,1],\ Mp=(r(p),\eta)\}.
\]
From the above formulation, it is easy to see that $G$ is maximally monotone if and only if there is no $p\in X^*$ such that
\[
r(p)\in[-1,1]
\qquad\text{and}\qquad
Mp=(r(p),\eta).
\]
This means that there is no $(p,\tau)\in X^*\times[-1,1]$ satisfying 
$r(p)=\tau$ and $Mp=(\tau,\eta)$, i.e.,  Eq. (\ref{eq:abs-endpoint-exclusion}).






\noindent\textit{\textup{(\ref{thm:sum-part})}}
Assume Eq. (\ref{eq:abs-endpoint-exclusion}) is true, by
part~\textup{(\ref{thm:maximality-part})}, the operator $\mathcal A$
with graph $G$ is maximally monotone.

For $a\in D$, Eq. (\ref{eq:abs-pairing}) and $Ma=F(r(a))$ give
\[
\begin{aligned}
|\pair{La}{g}|&=|r(a)|>1,\\
\pair{La}{a}
&=r(a)^2-\|Ma\|_H^2\\
&=r(a)^2-1-\|\omega(|r(a)|-1)\|_{H_0}^2>-S.
\end{aligned}
\]
Since
\[
|\pair{La}{g}|\le \|La\|\|g\|,
\]
we also have
\[
\|La\|>\frac{1}{\|g\|}.
\]
Therefore, $0\notin\dom\mathcal A$, and for every $a\in D$,
\[
\frac{\max\{0,-\pair{La}{a}\}}{\|La\|}
\le
\frac{S}{\|La\|}
<
S\|g\|.
\]
Taking the supremum gives
\[
\mathcal V(G)\le S\|g\|.
\]
Next, we define
\[
\mathcal Px=\pair{x}{g}g.
\]
The map $\mathcal P$ is bounded, linear, positive, and defined on all
of $X$. Since $g\ne0$ and
\[
\operatorname{ran}\mathcal P=\R g,
\]
it is nonzero and rank one.

Then, we prove that $\mathcal P$ is maximally monotone. Its monotonicity
follows from
\[
\pair{x-y}{\mathcal Px-\mathcal Py}
=
\pair{x-y}{g}^2\ge0.
\]
Let $(z,z^*)\in(\gra\mathcal P)^\mu$. Testing the polar inequality
against $(z+tv,\mathcal P(z+tv))\in\gra\mathcal P$ gives
\[
-t\pair{v}{z^*-\mathcal Pz}
+t^2\pair{v}{g}^2\ge0
\qquad(t\in\R,\ v\in X).
\]
Letting $t\to0$ through positive and negative values yields
\[
\pair{v}{z^*-\mathcal Pz}=0
\qquad(v\in X),
\]
which means $z^*=\mathcal Pz$. 
Thus $(\gra\mathcal P)^\mu=\gra\mathcal P$, which means $\mathcal P$ is maximally
monotone. 
Assumption~\ref{ass:construction}\textup{(\ref{ass:realization})}
with $t=2$ gives $a^2\in D$, which means 
$La^2\in\dom\mathcal A$. Since $\dom\mathcal P=X$, we have 
\[
\dom\mathcal A\cap\operatorname{int}(\dom\mathcal P)
=
\dom\mathcal A\ne\varnothing.
\]

Finally, for $a\in D$, Eq. (\ref{eq:abs-pairing}) gives
\[
\mathcal PLa
=
\pair{La}{g}g
=
r(a)g.
\]
Therefore, we infer 
\[
\gra(\mathcal A+\mathcal P)
=
\{(La,a+r(a)g):a\in D\}.
\]
For every such point,
\[
\begin{aligned}
\pair{La}{a+\mathcal PLa}
&=\pair{La}{a}+r(a)\pair{La}{g}\\
&=2r(a)^2-1-\|\omega(|r(a)|-1)\|_{H_0}^2\\
&>1-S>0.
\end{aligned}
\]
Therefore,
\[
(0,0)\in
\bigl(\gra(\mathcal A+\mathcal P)\bigr)^\mu.
\]
Moreover,
\[
\dom(\mathcal A+\mathcal P)
=
\dom\mathcal A,
\]
so $0\notin\dom(\mathcal A+\mathcal P)$ and consequently
\[
(0,0)\notin\gra(\mathcal A+\mathcal P).
\]
Thus
\[
(0,0)\in
\bigl(\gra(\mathcal A+\mathcal P)\bigr)^\mu
\setminus
\gra(\mathcal A+\mathcal P),
\]
and $\mathcal A+\mathcal P$ is not maximally monotone.
\end{proof}

\begin{lemma}[Pullback by a bounded surjection]
\label{lem:surjection-transport}
Let $U,V$ be real Banach spaces and let $Q:U\to V$ be a bounded linear surjection. Suppose $C_Q>0$ and every $y\in V$ has a preimage $u\in U$ satisfying
\[
\|u\|\le C_Q\|y\|.
\]
If $\mathcal M:V\rightrightarrows V^*$ is maximally monotone, then
\[
\gra(Q^*\mathcal M Q)
=
\{(u,Q^*a):(Qu,a)\in\gra\mathcal M\}
\]
is maximally monotone in $U\times U^*$, where $Q^*:V^*\to U^*$ is the adjoint of $Q$, and
\[
\pair{u}{Q^*a}=\pair{Qu}{a}
\qquad(u\in U,\ a\in V^*).
\]
\end{lemma}

\begin{proof}
Since $\mathcal M$ is maximally monotone, its graph is nonempty. Therefore, surjectivity of $Q$ implies that $\gra(Q^*\mathcal M Q)$ is nonempty. For any two points, we have 
$(u,Q^*a),(w,Q^*b)\in\gra(Q^*\mathcal M Q)$,
\[
\pair{u-w}{Q^*(a-b)}
=
\pair{Qu-Qw}{a-b}
\ge0.
\]
Thus $\gra(Q^*\mathcal M Q)$ is monotone.

Let $(v,v^*)\in U\times U^*$ be monotonically related to
$\gra(Q^*\mathcal M Q)$. Fix
$(u,Q^*a)\in\gra(Q^*\mathcal M Q)$. For every
$k\in\ker Q$ and $t\in\R$,
\[
(u+tk,Q^*a)\in\gra(Q^*\mathcal M Q),
\]
since $Q(u+tk)=Qu$. Hence
\[
\pair{v-u}{v^*-Q^*a}
-t\pair{k}{v^*}\ge0
\qquad(t\in\R).
\]
Since this holds for every $t\in\R$, we obtain
\[
\pair{k}{v^*}=0
\qquad(k\in\ker Q).
\]
Thus $v^*$ annihilates $\ker Q$.

For $y\in V$, choose any $w\in U$ satisfying $Qw=y$ and define
\[
p(y)=\pair{w}{v^*}.
\]
If $Qw_1=Qw_2=y$, then
$w_1-w_2\in\ker Q$, and hence
\[
\pair{w_1-w_2}{v^*}=0.
\]
The map $p$ is linear. Moreover, using a preimage $w$ satisfying
$\|w\|\le C_Q\|y\|$, we obtain
\[
|p(y)|
\le
\|w\|\|v^*\|
\le
C_Q\|v^*\|\|y\|.
\]
Thus $p\in V^*$ and, for every $w\in U$,
\[
\pair{w}{v^*}
=
p(Qw)
=
\pair{w}{Q^*p}.
\]
Therefore,
\[
v^*=Q^*p.
\]

Now let $(y,a)\in\gra\mathcal M$. By surjectivity, choose
$u\in U$ with $Qu=y$. Then
$(u,Q^*a)\in\gra(Q^*\mathcal M Q)$, so the polar inequality gives
\[
\pair{v-u}{v^*-Q^*a}\ge0.
\]
Since $v^*=Q^*p$,
\[
\pair{Qv-y}{p-a}\ge0
\qquad((y,a)\in\gra\mathcal M).
\]
Thus, $(Qv,p)$ is monotonically related to $\gra\mathcal M$. By maximality of $\mathcal M$,
\[
(Qv,p)\in\gra\mathcal M.
\]
This means
\[
(v,Q^*p)=(v,v^*)\in\gra(Q^*\mathcal M Q).
\]
Thus every point monotonically related to
$\gra(Q^*\mathcal M Q)$ already belongs to the graph. Since the graph is monotone, $\gra(Q^*\mathcal M Q)$ is maximally monotone. 
\end{proof}

\section{Counterexamples on $c_0$ and standard $\ell^1$}
\label{sec:counterexamples}
In this section, we apply Theorem~\ref{thm:endpoint-criterion} to construct a counterexample on $c_0$ and then use Lemma~\ref{lem:surjection-transport} to obtain a counterexample on standard $\ell^1$.


% We first verify the structural theorem for explicit operators on $c_0$.
% A bounded surjection then defines operators on standard $\ell^1$, whose
% maximality and sum failure are proved in the full continuous dual pair.

\subsection{A counterexample on $c_0$}

Let $I=\Nzero\times\N$ and $Y=c_0(I)$ with the supremum norm. For every $x\in Y$ and $\epsilon>0$, only finitely many $(b,j)\in I$ satisfy $|x_{b,j}|\ge\epsilon$. Its continuous dual is $\ell^1(I)$, with $\pair{x}{a}=\sum_{b,j}x_{b,j}a_{b,j}$. We also use this pairing for $x\in\ell^\infty(I)$ and $a\in\ell^1(I)$. The unit coordinate vectors are denoted by $e_{b,j}$.

We index coordinates by blocks. The map $m$ records the sum in each block, and $E$ applies a triangular operator within that block. The additional constraint below couples the block sums through $F$.
\begin{definition}[Block maps]
For $a\in\ell^1(I)$, define
\begin{equation}
\begin{aligned}
(ma)_b&=\sum_{j\ge1}a_{b,j},\\
(Ea)_{b,j}&=a_{b,j}+2\sum_{k>j}a_{b,k}.
\end{aligned}
\label{eq:D4}
\end{equation}
Here $m:\ell^1(I)\to\ell^1(\Nzero)\subset\ell^2(\Nzero)$ and $E:\ell^1(I)\to Y$. Put
\begin{equation}
\begin{aligned}
g&=e_{0,1}-e_{0,2}\in Y^*,\\
h&=Eg=-e_{0,1}-e_{0,2}\in Y,\\
r(a)&=\pair{h}{a}=-a_{0,1}-a_{0,2},\\
La&=-Ea+r(a)h.
\end{aligned}
\label{eq:D5}
\end{equation}
\end{definition}

\begin{definition}[The curve of block sums]
For $s>0$ and $n\ge1$, write
\begin{equation}
\begin{aligned}
\omega_s(n)&=\frac{\max\{1-sn^{\frac{1}{4}},0\}}{4n},\\
W(s)&=\sum_{n\ge1}\omega_s(n)^2,\\
\sigma&=\sum_{n\ge1}\frac1{16n^2}=\frac{\pi^2}{96}.
\end{aligned}
\label{eq:D6}
\end{equation}
Each $\omega_s$, $s>0$, has finite support. Set $\omega_0=(\frac{1}{4n})_{n\ge1}\in\ell^2\setminus\ell^1$. Define
\begin{equation}
\begin{aligned}
J&=(-\infty,-1)\cup(1,\infty),\\
F(t)&=(\sgn t,\omega_{|t|-1})\in\ell^1(\Nzero)\quad(t\in J),\\
D&=\{a\in\ell^1(I):r(a)\in J,\ ma=F(r(a))\}.
\end{aligned}
\label{eq:D7}
\end{equation}

\end{definition}

\begin{definition}[The two operators on $c_0(I)$]
The operators and kernel are
\begin{equation}
\begin{aligned}
\gra A&=\{(La,a):a\in D\},\\
K&=\{k\in\ell^1(I):mk=0,\ r(k)=0\},\\
Px&=\pair{x}{g}g\quad(x\in Y).
\end{aligned}
\label{eq:D8}
\end{equation}
For each $t\in J$, a vector satisfying $r(a^t)=t$ and $ma^t=F(t)$ is
\begin{equation}
a^t=-t e_{0,1}+(t+\sgn t)e_{0,3}
       +\sum_{n\ge1}\omega_{|t|-1}(n)e_{n,1}.
\label{eq:D9}
\end{equation}
\end{definition}

% Coordinate verification follows the definitions in the first example.
\label{sec:supporting}





We now verify the hypotheses of Theorem~\ref{thm:endpoint-criterion} for the operators in Eq. (\ref{eq:D8}). Lemmas~\ref{lem:coordinate-pairing} and~\ref{lem:coordinate-annihilator} establish \textup{(\ref{ass:bilinear})} and \textup{(\ref{ass:annihilator})} of Assumption~\ref{ass:construction}. Lemma~\ref{lem:coordinate-profile} verifies the curve and fibre conditions \textup{(\ref{ass:curve-limit})}--\textup{(\ref{ass:realization})} and proves the exclusion condition in Eq. (\ref{eq:abs-endpoint-exclusion}).

\begin{lemma}[The coordinate pairing identity]
\label{lem:coordinate-pairing}
The maps $m:\ell^1(I)\to\ell^1(\Nzero)\subset\ell^2(\Nzero)$ and $E,L:\ell^1(I)\to c_0(I)$ are bounded. For all $a,c\in\ell^1(I)$, Eqs. (\ref{eq:C1}) and (\ref{eq:C2}) below hold.
\end{lemma}
\begin{proof}
\label{sec:C1}
For $a\in\ell^1(I)$, $\|ma\|_1\le\|a\|_1$ and $\|Ea\|_\infty\le2\|a\|_1$. If $a$ has finite support, then $Ea$ also has finite support. Indeed, only finitely many blocks contain nonzero coordinates, and within each such block no output occurs after the last nonzero coordinate of $a$. Since finitely supported vectors are dense in $\ell^1(I)$, the bound on $E$ then implies $Ea\in c_0(I)$ for every $a\in\ell^1(I)$. Thus $L$ also maps $\ell^1(I)$ boundedly into $c_0(I)$. 

For $a,c\in\ell^1(I)$, absolute convergence permits exchanging the double sums. In each block the diagonal terms and the two triangular off-diagonal sums give
\begin{equation}
\pair{Ea}{c}+\pair{Ec}{a}
 =2\sum_b\Big(\sum_j a_{b,j}\Big)\Big(\sum_j c_{b,j}\Big)
 =2(ma,mc)_{\ell^2}.
\label{eq:C1}
\end{equation}
The absolute sum of all products is at most $\|a\|_1\|c\|_1$ before the factor $2$, so the identity holds for all $a,c\in\ell^1(I)$. Since $mg=0$, $h=Eg$, and $r(g)=0$, Eq. (\ref{eq:C1}) implies $\pair{Ea}{g}=-r(a)$. Consequently
\begin{equation}
\begin{aligned}
\pair{La}{a}&=r(a)^2-\|ma\|_2^2,\\
\pair{La}{g}&=r(a),\\
\pair{La-Lc}{a-c}&=(r(a)-r(c))^2-\|ma-mc\|_2^2.
\end{aligned}
\label{eq:C2}
\end{equation}
\end{proof}

\begin{lemma}[The annihilator of $K$ in $c_0(I)$]
\label{lem:coordinate-annihilator}
For $K$ in Eq. (\ref{eq:D8}),
\[
\{z\in c_0(I):\pair{z}{k}=0\ \text{for every }k\in K\}=\R h.
\]
\end{lemma}
\begin{proof}
Let $z\in c_0(I)$ annihilate $K$. In a block $b\ge1$, every difference $e_{b,j}-e_{b,k}$ is in $K$. Thus $z$ is constant in that entire block. Since $z\in c_0(I)$, this constant must be zero. In block zero, differences between indices $j,k\ge3$ belong to $K$, so the common tail value is also zero. Finally, $g=e_{0,1}-e_{0,2}$ belongs to $K$, forcing $z_{0,1}=z_{0,2}$. It follows that $z$ is a scalar multiple of $h$. Conversely, $\pair{h}{k}=r(k)=0$ proves that every multiple of $h$ annihilates $K$.
\end{proof}

\begin{lemma}[Properties of the curve and its parameter fibres]
\label{lem:coordinate-profile}
The map $s\mapsto\omega_s$ from $(0,\infty)$ to $\ell^2(\N)$ satisfies Eq. (\ref{eq:abs-tail}) with $\eta=\omega_0$, $S=\sigma<\frac{1}{8}$ and a Lipschitz constant strictly less than one. For every $t\in J$, the vector $a^t$ in Eq. (\ref{eq:D9}) belongs to $D$, and its parameter fibre is $a^t+K$. No $p\in\ell^1(I)$ and $\tau\in[-1,1]$ satisfy $r(p)=\tau$ and $mp=(\tau,\omega_0)$.
\end{lemma}
\begin{proof}
\label{sec:C2}
For each $s>0$, $\omega_s(n)=0$ whenever $n\ge s^{-4}$. Also $0\le\omega_s(n)\le\frac{1}{4n}$. Thus, $W(s)\le\sigma<\frac{1}{8}$, and dominated convergence in $\ell^2$ gives $\omega_s\to\omega_0$ and $W(s)\to\sigma$ as $s$ decreases to zero. The comparison vector $\omega_0$ is not in $\ell^1$ by divergence of the harmonic series.

Since $v\mapsto\max\{v,0\}$ is $1$-Lipschitz on $\R$,
\begin{equation}
\|\omega_s-\omega_q\|_2^2\le c|s-q|^2,\qquad
c=\frac1{16}\sum_{n\ge1}n^{-\frac{3}{2}}<\frac3{16}<1.
\label{eq:C3}
\end{equation}
The strict numerical bound follows from $\sum_{n\ge1}n^{-\frac{3}{2}}<1+\int_1^\infty x^{-\frac{3}{2}}\,dx=3$. For $t,u$ on the same branch of $J$, Eq. (\ref{eq:C3}) gives $\|F(t)-F(u)\|_2\le|t-u|$. For $t=1+s$ and $u=-1-q$, the squared distance is at most $4+c(s-q)^2\le(2+s+q)^2$, and the opposite orientation is identical. Thus $F$ is $1$-Lipschitz on all of $J$.

The finitely supported vector $a^t$ in Eq. (\ref{eq:D9}) satisfies $r(a^t)=t$ and $ma^t=F(t)$, proving nonemptiness for every $t\in J$. The parameter fibre $\{a:r(a)=t,\ ma=F(t)\}$ is precisely $a^t+K$. For every $p\in\ell^1(I)$, the vector $mp$ is in $\ell^1(\Nzero)$. Since $\omega_0\notin\ell^1(\N)$, the last assertion follows.
\end{proof}

% The first counterexample follows its coordinate lemmas.

The preceding lemmas establish all the hypotheses of Theorem~\ref{thm:endpoint-criterion}, including Eq. (\ref{eq:abs-endpoint-exclusion}). Therefore, applying that theorem to the operators in Eq. (\ref{eq:D8}) gives the following counterexample. 

\begin{theorem}[A counterexample on $c_0$]\label{thm:C}
Let $I=\Nzero\times\N$ and $Y=c_0(I)$ with its usual supremum norm and continuous dual $Y^*=\ell^1(I)$. Let $A:Y\rightrightarrows Y^*$ and the bounded positive rank-one map $P:Y\to Y^*$ be defined by Eq. (\ref{eq:D8}). Then 
\begin{enumerate}
\item\label{thm:C-maximality} $A$ and $P$ are maximally monotone in $Y\times Y^*$.
\item\label{thm:C-domain} $\dom P=Y$ and $\dom A\ne\varnothing$, so
$\dom A\cap\operatorname{int}(\dom P)\ne\varnothing$.
\item\label{thm:C-bound} Every $(x,a)\in\gra A$ satisfies $\|x\|_\infty>\frac{1}{2}$ and
$\pair{x}{a}\ge-2\sigma\|x\|_\infty$, where $\sigma=\frac{\pi^2}{96}<\frac{1}{8}$.
\item\label{thm:C-pairing} Every $(x,b)\in\gra(A+P)$ satisfies
$\pair{x}{b}>1-\sigma>\frac{7}{8}$.
\item\label{thm:C-sum} $(0,0)$ is monotonically related to the entire graph of $A+P$ and
does not belong to that graph. Thus $A+P$ is not maximally monotone.
\end{enumerate}
\end{theorem}
\begin{proof}
\noindent\textit{\textup{(\ref{thm:C-maximality})} and \textup{(\ref{thm:C-domain})}. Maximality and the domain intersection condition.}
Apply Theorem~\ref{thm:endpoint-criterion} with $X=Y$, $M=m$, $H_0=\ell^2(\N)$, $\omega(s)=\omega_s$, $\eta=\omega_0$ and $S=\sigma$. Lemmas~\ref{lem:coordinate-pairing}--\ref{lem:coordinate-profile}, together with $g\ne0$ and $mg=0$ from Eq. (\ref{eq:D5}), verify Assumption~\ref{ass:construction}. The last assertion of Lemma~\ref{lem:coordinate-profile} gives Eq. (\ref{eq:abs-endpoint-exclusion}), so the theorem proves maximal monotonicity of both $A$ and $P$. 

\label{sec:C6}
Since $\|g\|_1=2$, we have $\|P\|\le4$ and $Pe_{0,1}=g\ne0$. The vector $a^2=-2e_{0,1}+3e_{0,3}$ in Eq. (\ref{eq:D9}) belongs to $D$ and gives $La^2=-6e_{0,1}-8e_{0,2}-3e_{0,3}$. This means $La^2\in\dom A$, and $\dom P=Y$ yields $\dom A\cap\operatorname{int}(\dom P)\ne\varnothing$.

\noindent\textit{\textup{(\ref{thm:C-bound})}. Bounds on $\gra A$.}
\label{sec:C5}
For every $a\in D$, $t=r(a)$ satisfies $|t|>1$. By Eq. (\ref{eq:C2}) and $\|g\|_1=2$,
\begin{equation}
\begin{aligned}
\|La\|_\infty&\ge\frac{|t|}{2}>\frac{1}{2},\\
\pair{La}{a}&=t^2-1-W(|t|-1)>-\sigma.
\end{aligned}
\label{eq:C9}
\end{equation}
Eq. (\ref{eq:C9}) gives $A(0)=\varnothing$ and, for every $a\in D$,
\[
\frac{\max\{0,-\pair{La}{a}\}}{\|La\|_\infty}
\le \frac{\sigma}{\|La\|_\infty}
<2\sigma.
\]
Therefore $\pair{La}{a}\ge-2\sigma\|La\|_\infty$ for every $a\in D$, and taking the supremum gives $\mathcal V(\gra A)\le2\sigma<\frac{1}{4}$.

\noindent\textit{\textup{(\ref{thm:C-pairing})} and \textup{(\ref{thm:C-sum})}. Nonmaximality of $A+P$.}
To verify nonmaximality of $A+P$, take an arbitrary point of $\gra(A+P)$. By Eqs. (\ref{eq:D8}) and (\ref{eq:C2}), it has the form $(La,a+tg)$ with $a\in D$ and $t=r(a)$. Thus Eq. (\ref{eq:C2}) gives
\begin{equation}
\pair{La}{a+tg}=2t^2-1-W(|t|-1)>1-\sigma>\frac{7}{8}.
\label{eq:C11}
\end{equation}
Eq. (\ref{eq:C11}) shows that $(0,0)$ is monotonically related to every point of $\gra(A+P)$. Since $A$ and $P$ are monotone, $\gra(A+P)\cup\{(0,0)\}$ is monotone. Eq. (\ref{eq:C9}) gives $0\notin\dom(A+P)$, so this union strictly contains $\gra(A+P)$. Therefore $A+P$ is not maximally monotone.
\end{proof}

% Consequences of the first counterexample.
\label{sec:domain-comparison}

% \begin{corollary}\label{prop:domain-gap}
% For the operator $A$ of Example~\ref{thm:C},
% $0\in\operatorname{conv}(\dom A)$ and
% $\operatorname{dist}(0,\dom A)\ge1/2$.
% In particular, $\overline{\dom A}$ is not convex.
% \end{corollary}
% \begin{proof}
% At $t=2$ and $t=-2$, the tails in Eq. (\ref{eq:D9}) vanish, so that
% $a^{-2}=-a^2$. Since $L$ is linear, both $La^2$ and $-La^2$ belong
% to $\dom A$. Their midpoint is zero. Eq. (\ref{eq:C9}) gives
% $\|x\|_\infty>1/2$ for every $x\in\dom A$, and hence the asserted
% distance bound. A convex closure containing $La^2$ and $-La^2$ would
% contain zero, contradicting this bound.
% \end{proof}



% \begin{remark}[Boundary points and unbounded dual values]
% \label{sec:structure}
% The nonsummable limiting block sums are also visible along a sequence
% of points in $\gra A$ whose primal components converge in norm.
% Set $t_n^\pm=\pm(1+1/n)$, $a_n^\pm=a^{t_n^\pm}$, and
% $x_n^\pm=La_n^\pm$. For $t\in J$ and $\epsilon=\sgn t$,
% direct substitution in Eqs. (\ref{eq:D4})--(\ref{eq:D9}) gives
% \[
% (La^t)_{0,1}=-2t-2\epsilon,\qquad
% (La^t)_{0,2}=-3t-2\epsilon,\qquad
% (La^t)_{0,3}=-t-\epsilon,
% \]
% with the remaining block-zero coordinates zero. In block $b\ge1$,
% the only nonzero coordinate is $-\omega_{|t|-1}(b)$ at index $1$.
% Consequently, $x_n^\pm$ converge in the supremum norm to
% $x_\pm\in c_0(I)$ with block-zero coordinates $\mp(4,5,2)$ and
% positive-block first coordinates $-1/(4b)$. Indeed, convergence of
% the tails in $\ell^2$ implies convergence in the supremum norm.
% The limits satisfy $\pair{x_\pm}{g}=\pm1$ and lie outside $\dom A$.
% On the other hand,
% \[
% \|a_n^\pm\|_1=2(1+1/n)+1+\|\omega_{1/n}\|_1\longrightarrow\infty.
% \]
% For each fixed $N$, the sum of the first $N$ tail coordinates tends to
% $\sum_{b=1}^N1/(4b)$; these partial sums are unbounded. This proves
% the divergence without asserting convergence of the labels in the dual.

% \end{remark}

\subsection{A counterexample on standard $\ell^1$}

Next, we transfer the counterexample of Theorem~\ref{thm:C} from $Y=c_0(I)$ to standard $\ell^1$ through a bounded linear surjection.
% Lemma~\ref{lem:surjection-transport} will establish the maximality of the transferred operators, and the adjoint identity will preserve the pairing inequality that proves nonmaximality of their sum.

% The following construction makes explicit the standard representation of a separable Banach space as a quotient of $\ell^1$ \cite[Section~1.12]{megginson1998}.

Let $Z=\ell^1(\N)$ and $Z^*=\ell^\infty(\N)$, with their usual pairing. To construct the required surjection, fix an enumeration $(q_n)$ of all finitely supported rational vectors in the closed unit ball of $Y$.
\begin{definition}[The operators on standard $\ell^1$]
Define
\begin{equation}
\begin{aligned}
Qu&=\sum_{n\ge1}u_nq_n,\\
(Q^*a)_n&=\pair{q_n}{a},\\
\gra T&=\{(u,Q^*a):a\in D,\ Qu=La\},\\
f&=Q^*g,\\
Bu&=\pair{u}{f}f.
\end{aligned}
\label{eq:D10}
\end{equation}
\end{definition}
The next lemma provides the bounded preimages needed in Lemma~\ref{lem:surjection-transport}, and no bounded linear right inverse of $Q$ is required.

\begin{lemma}[The specified quotient onto $c_0$]
\label{lem:specified-quotient}
The map $Q:Z=\ell^1(\N)\to Y=c_0(I)$ of Eq. (\ref{eq:D10}) is a bounded surjection with $\|Q\|=1$. Its adjoint satisfies $\|Q^*a\|_\infty=\|a\|_1$ for every $a\in\ell^1(I)$. For the fixed enumeration $(q_n)$, a map $\mathcal R:Y\to Z$ can be specified with
\[
Q\mathcal R(x)=x,\qquad \|\mathcal R(x)\|_1\le2\|x\|_\infty.
\]
\end{lemma}
\begin{proof}
\label{sec:L1}
Finitely supported rational vectors in the closed unit ball of $Y$ are countable and dense in that ball: truncate a $c_0$ vector and approximate each retained coordinate by a rational inside $[-1,1]$. Fix an enumeration of all these vectors as $(q_n)$. For $u\in Z=\ell^1(\N)$, the sum defining $Qu$ converges absolutely in the Banach space $Y$, and $\|Qu\|_\infty\le\|u\|_1$. Coordinate unit vectors occur among the $q_n$, so $\|Q\|=1$.

For an arbitrary $x\in Y$, set $r_0=x$. If $r_k\ne0$, choose the least index $n_k$ for which $q_{n_k}$ satisfies $\|q_{n_k}-\frac{r_k}{\|r_k\|_\infty}\|_\infty<\frac{1}{2}$ and set
\begin{equation}
c_k=\|r_k\|_\infty,\qquad r_{k+1}=r_k-c_kq_{n_k}.
\label{eq:L1}
\end{equation}
If a residual vanishes, stop. Otherwise $\|r_k\|_\infty\le2^{-k}\|x\|_\infty$ and $\sum_k c_k\le2\|x\|_\infty$. Define $u_n=\sum_{k:n_k=n}c_k$. Repeated indices are combined. Positivity of the $c_k$ and summability give $u\in\ell^1$ with $\|u\|_1=\sum_k c_k\le2\|x\|_\infty$. Telescoping Eq. (\ref{eq:L1}), then absolute convergence under regrouping, proves
\begin{equation}
Qu=x,\qquad \|u\|_1\le2\|x\|_\infty.
\label{eq:L2}
\end{equation}
For $x=0$, use $u=0$. This proves surjectivity with the asserted lifting bound.

\label{sec:L2}
Absolute convergence gives $\pair{u}{Q^*a}=\pair{Qu}{a}$ and $\|Q^*a\|_\infty\le\|a\|_1$. For each finite subset of $I$, its sign vector for $a$ is a finitely supported rational vector of norm at most one and occurs in the enumeration. Taking these finite sign tests proves
\begin{equation}
\|Q^*a\|_\infty=\|a\|_1\qquad(a\in\ell^1(I)).
\label{eq:L3}
\end{equation}
In particular, $Q^*$ is injective. The least-index choices in Eq. (\ref{eq:L1}) specify
\[
\mathcal R(x)=\sum_k c_k e_{n_k},\qquad \mathcal R(0)=0.
\]
The series converges in $\ell^1$ by the bound on $\sum_k c_k$, and Eq. (\ref{eq:L2}) gives both asserted properties.
\end{proof}

By Lemma~\ref{lem:specified-quotient}, $Q\mathcal R(La)=La$, so the solutions of $Qu=La$ are exactly $u=\mathcal R(La)+v$ with $v\in\ker Q$. Substituting this expression into Eq. (\ref{eq:D10}) gives
\begin{equation}
\gra T=\{(\mathcal R(La)+v,Q^*a):a\in D,\ v\in\ker Q\}.
\label{eq:L4}
\end{equation}
This formula includes every preimage of each point in $\dom A$. We can therefore apply Lemma~\ref{lem:surjection-transport} to the counterexample in Theorem~\ref{thm:C}. Together with the pairing identities below, this yields the following counterexample on standard $\ell^1$.


\begin{theorem}[A counterexample on standard $\ell^1$]\label{thm:L}
On $Z=\ell^1(\N)$ with its usual norm and continuous dual $Z^*=\ell^\infty(\N)$, let $T:Z\rightrightarrows Z^*$ and $B:Z\to Z^*$ be defined by Eq. (\ref{eq:D10}). Then:
\begin{enumerate}
\item\label{thm:L-maximality} $T$ is maximally monotone in the full dual pair $Z\times Z^*$.
\item\label{thm:L-partner} $B$ is nonzero, bounded, positive, rank-one, and maximally monotone,
with $\dom B=Z$. In particular, $\dom T\cap\operatorname{int}(\dom B)\ne\varnothing$.
\item\label{thm:L-bound} $T(0)=\varnothing$, and every $(u,b)\in\gra T$ satisfies
$\|u\|_1>\frac{1}{2}$ and $\pair{u}{b}\ge-2\sigma\|u\|_1$. Thus $\mathcal V(\gra T)\le2\sigma<\frac{1}{4}$.
\item\label{thm:L-pairing} Every $(u,c)\in\gra(T+B)$ satisfies
$\pair{u}{c}>1-\sigma>\frac{7}{8}$.
\item\label{thm:L-sum} $(0,0)$ belongs to the monotone polar of $\gra(T+B)$ and does
not belong to $\gra(T+B)$. Hence $T+B$ is not maximally monotone.
\end{enumerate}
\end{theorem}
\begin{proof}
\noindent\textit{\textup{(\ref{thm:L-maximality})} and \textup{(\ref{thm:L-partner})}. Maximality and the domain intersection condition.}
Eq. (\ref{eq:L4}) identifies $T$ with $Q^*AQ$, including every preimage under $Q$. Lemma~\ref{lem:specified-quotient} supplies the preimage bound with $C_Q=2$. Since $A$ is maximally monotone by Theorem~\ref{thm:C}, Lemma~\ref{lem:surjection-transport} proves maximal monotonicity of $T$ in $Z\times Z^*=\ell^1\times\ell^\infty$.

\label{sec:L4}
Let $f=Q^*g$. By Eq. (\ref{eq:L3}), $\|f\|_\infty=\|g\|_1=2$, so $f$ is nonzero. The map $B:u\mapsto\pair{u}{f}f$ is nonzero, linear, bounded with norm at most $4$, positive, and rank-one. Moreover, for every $u\in Z$,
\[
Bu=\pair{u}{Q^*g}Q^*g
   =Q^*(\pair{Qu}{g}g)
   =Q^*(P(Qu)).
\]
Thus $B=Q^*PQ$. Maximal monotonicity of $P$ in Theorem~\ref{thm:C} and Lemma~\ref{lem:surjection-transport} give maximal monotonicity of $B$. Its domain is all of $Z$, so $\operatorname{int}(\dom B)=Z$.

The vector $q_\circ=-(\frac{3}{4})e_{0,1}-e_{0,2}-(\frac{3}{8})e_{0,3}$ occurs as $q_{n_0}$. For $u_\circ=8e_{n_0}$, $Qu_\circ=La^2$. Thus $(u_\circ,Q^*a^2)\in\gra T$, and $\dom T\cap\operatorname{int}(\dom B)\ne\varnothing$.

\noindent\textit{\textup{(\ref{thm:L-bound})}. Bounds on $\gra T$.}
\label{sec:L3}
For $(u,Q^*a)\in\gra T$, we have $Qu=La$. Thus the adjoint identity, $\|Q\|=1$ and Eq. (\ref{eq:C9}) transfer the bounds for $A$ to $T$. Writing $t=r(a)$, we obtain
\begin{equation}
\begin{aligned}
\pair{u}{Q^*a}&=\pair{La}{a}
 =t^2-1-W(|t|-1)>-\sigma,\\
\|u\|_1&\ge\|Qu\|_\infty=\|La\|_\infty>\frac{1}{2}.
\end{aligned}
\label{eq:L6}
\end{equation}
Thus $T(0)=\varnothing$ and $\pair{u}{Q^*a}\ge-2\sigma\|u\|_1$ on the entire graph, including every $\ker Q$ translate. Consequently, $\mathcal V(\gra T)\le2\sigma<\frac{1}{4}$.

\noindent\textit{\textup{(\ref{thm:L-pairing})} and \textup{(\ref{thm:L-sum})}. Nonmaximality of $T+B$.}
\label{sec:L5}
For $(u,Q^*a)\in\gra T$, we have $Qu=La$ and $Bu=Q^*(P(La))$. The adjoint identity therefore transfers Eq. (\ref{eq:C11}) to $T+B$. With $t=r(a)$,
\begin{equation}
\begin{aligned}
\pair{u}{Q^*a+Bu}
 &=\pair{La}{a+P(La)}\\
 &=2t^2-1-W(|t|-1)>1-\sigma>\frac{7}{8}.
\end{aligned}
\label{eq:L7}
\end{equation}
Eq. (\ref{eq:L7}) shows that $(0,0)$ is monotonically related to every point of $\gra(T+B)$. Since $T$ and $B$ are monotone, $\gra(T+B)\cup\{(0,0)\}$ is monotone. Since $T(0)=\varnothing$, this union strictly contains $\gra(T+B)$. Therefore, $T+B$ is not maximally monotone.
\end{proof}

\section{Conclusions}
\label{sec:conclusion}

In this paper, we constructed counterexamples to Rockafellar's sum conjecture. Specifically, we established a general construction theorem that computes the monotone polar of the proposed graphs, characterizes their maximal monotonicity and identifies when adding an everywhere-defined positive rank-one operator produces a nonmaximal sum. We verified its hypotheses for explicit operators on $c_0$ and transferred the resulting pair to standard $\ell^1$ through a bounded linear surjection. In both pairs, the two operators are maximally monotone and satisfy the original interior-domain condition, and their sum is not maximally monotone. We also proved a finite radial bound for the first operator in each pair.

\section*{Code availability.}
The Lean formalization of the counterexample on $c_0$ and the general pullback lemma is available at \url{https://github.com/Weifeng-Yang/RockafellarSumConjecture}


\section*{Use of AI}
The author supplied the prior results and unsuccessful approaches, including some constructions based on blockwise triangular operators, together with a framework for positive rank-one perturbations, and guided their further development to investigate how maximality can fail under addition. The author also supplied obstruction results showing that a finite radial bound on a monotone graph need not survive maximal extension. Through iterative discussions of these materials, the author and GPT-5.6 Sol developed a general construction theorem characterizing maximality within this framework and identifying when an everywhere-defined rank-one monotone perturbation yields a nonmaximal sum.  Under the author's direction and using the aforementioned results as inputs, GPT-5.6 Sol worked out the detailed constructions and proof arguments, and the transfer from $c_0$ to standard $\ell^1$. 




% and directed the counterexample search and a general construction theorem characterizing maximality and showing when addition of an everywhere-defined rank-one monotone operator produces a nonmaximal sum, and. Under this direction, GPT-5.6 Sol developed explicit formulas and proof arguments for the counterexamples on $c_0$ and standard $\ell^1$, and assisted with literature checks, adversarial review and exposition. 

\begingroup
\raggedright \bibliographystyle{IEEEtran}
\bibliography{refer}
\endgroup

\end{document}
